\chapter{About this book}

Hi reader! o/

This mini-book is an adaptation of the Multisig for
Protocols framework from the Security Frameworks, an
initiative by the Security Alliance. It is written for
the person who is about to sign, or who is designing
the set those people will live in.

A threshold is not a security model. Cryptography
already works. Protocol multisigs fail operationally:
wrong sizing, one God-mode Safe, skipped verification,
a signer who is actually three laptops in the same
bag.

This is not, by any means, an extensive guide or
expected to be followed to the T. Its intention is to
guide and provide hints as to where to apply security.
Your threat model still wins.

You can access the contents of this book, and a
possibly more updated version, on the Security
Frameworks website.

Feel free to suggest an update by contributing
directly there.

\section{How to use this book}

Three words. If the rest falls out of the bag, keep
these:

\begin{enumerate}
  \item \textbf{Separate.} Upgrade, pause, and treasury
  must not live on one set.
  \item \textbf{Delay.} A timelock without someone
  watching is decoration.
  \item \textbf{Verify.} You authorize an effect, not
  a hash. Decode it. Safe is the default stack here
  because that is what most protocol sets actually
  run. The checks travel. The buttons do not.
\end{enumerate}

If you are joining a set you did not design, read
``Who may sign'' and ``How you join'' before you
accept. You can refuse a bad set.

\keepblock{%
If you value these type of resources, please consider
supporting them. You can do so by going to our website.

{\centering\plainurl{securityalliance.org/donate}\par}
}
